% Source: https://artofproblemsolving.com/community/c3h1200332_morbid_paradox
% Stable identity: aops-topic-1200332-q01
The products the student computes are the probabilities of the events ``I survive the first $n-1$ seconds and die in second $n$''. These events are disjoint, so their probabilities must add up to at most $1$; of course they get smaller as $n$ grows. That says nothing about the risk at a later time.

The relevant quantity is the conditional probability of dying in second $n$ \emph{given that I am alive at its start}, which is about $\dfrac{2}{N+2(n-1)}$ each time, essentially constant. With a roughly constant risk per second, the probability of surviving $n$ seconds tends to $0$, so everyone eventually dies; there is no paradox. (In reality the risk also rises with age, which this model ignores.)
